\section{Modelling and linearization}

\subsection{Ctrl+F keywords}

\subsubsection{RLC, impedance, mass spring damper, equations of motion, modellering}
\[
\text{Newton, force balance, voltage divider, white-box model, linearization, operating point}
\]
Use this section when the exam asks for a model before doing control design.

\subsection{Mechanical modelling}

\subsubsection{Newton's second law}
\[
\sum F=m\ddot{x}
\]
Choose a positive direction and keep spring, damper, and external-force signs consistent.

\subsubsection{Spring and damper forces}
\[
F_k=k(x_1-x_2),\qquad F_b=b(\dot{x}_1-\dot{x}_2)
\]
For coupled masses, relative displacement and relative velocity determine internal forces.

\subsubsection{Mass-spring-damper transfer function}
\[
m\ddot{x}+b\dot{x}+kx=F
\quad\Longrightarrow\quad
\frac{X(s)}{F(s)}=\frac{1}{ms^2+bs+k}
\]
With zero initial conditions, force-balance equations become algebraic transfer functions.

\subsection{Electrical modelling}

\subsubsection{Laplace impedances}
\[
Z_R=R,\qquad Z_L=sL,\qquad Z_C=\frac{1}{sC}
\]
Use impedances for RLC circuit transfer functions in the \(s\)-domain.

\subsubsection{Voltage divider}
\[
\frac{V_2(s)}{V_{\mathrm{in}}(s)}=\frac{Z_2(s)}{Z_1(s)+Z_2(s)}
\]
Pick \(Z_2\) as the impedance across the measured output voltage.

\subsection{Linearization}

\subsubsection{Equilibrium point}
\[
0=f(x_0,u_0),\qquad y_0=g(x_0,u_0)
\]
The operating point must satisfy the nonlinear state equation for constant input.

\subsubsection{First-order Taylor linearization}
\[
\Delta\dot{x}=A\Delta x+B\Delta u,\qquad
\Delta y=C\Delta x+D\Delta u
\]
\[
A=\left.\frac{\partial f}{\partial x}\right|_0,\quad
B=\left.\frac{\partial f}{\partial u}\right|_0,\quad
C=\left.\frac{\partial g}{\partial x}\right|_0,\quad
D=\left.\frac{\partial g}{\partial u}\right|_0
\]
The variables are deviations from the operating point.

\subsubsection{State-space to transfer function}
\[
G(s)=C(sI-A)^{-1}B+D
\]
Use this only after selecting the correct input and output channels.

\subsection{Exam recipe: RLC circuit transfer function}

\textbf{Use when:} the problem shows resistors, inductors, capacitors, and asks for \(V_{\mathrm{out}}/V_{\mathrm{in}}\) or a block diagram.

\textbf{Steps:}
\begin{enumerate}
\item Replace \(R,L,C\) by \(R,sL,1/(sC)\).
\item Identify series and parallel equivalent impedances.
\item Use the voltage divider or Kirchhoff equations.
\item Simplify to a rational function in \(s\).
\item Check low- and high-frequency limits.
\end{enumerate}

\textbf{Fast checks:} capacitors are open at DC and short at high frequency; inductors are short at DC and open at high frequency.

\textbf{Common traps:} choosing the wrong output impedance in the voltage divider.

\subsection{Exam recipe: coupled mass equations of motion}

\textbf{Use when:} the problem asks for equations of motion or transfer function of masses, springs, and dampers.

\textbf{Steps:}
\begin{enumerate}
\item Draw one free-body equation for each mass.
\item Use relative displacement for springs and relative velocity for dampers.
\item Apply \(\sum F=m\ddot{x}\) to each mass.
\item Laplace transform with zero initial conditions.
\item Solve the algebraic equations for the requested output/input ratio.
\end{enumerate}

\textbf{Fast checks:} internal spring and damper forces appear with opposite signs on the two connected masses.
